\capitulo{3}{Conceptos Teóricos}

A continuación, se detallan los principios teóricos que sustentan la arquitectura del sistema.

\section{Contexto sanitario}

\textbf{Enfermedad de Parkinson.} Se define como un trastorno neurodegenerativo crónico y progresivo que afecta al sistema nervioso central, específicamente a las neuronas encargadas de producir dopamina. Principalmente, esta enfermedad se manifiesta en personas de edad avanzada, aumentando su incidencia significativamente a partir de los 60 años. Desde la perspectiva de la interacción persona-computador, esta patología se caracteriza por presentar síntomas motores severos como temblor en reposo, rigidez muscular y bradicinesia. La comprensión teórica de esta sintomatología es el requisito previo indispensable para el diseño de la interfaz de usuario, imponiendo restricciones estrictas de accesibilidad para garantizar que el software sea operable por personas con control motor fino reducido. \cite{mayo_parkinson}

\textbf{Bradicinesia.} Es el término clínico que describe la lentitud en la ejecución de los movimientos voluntarios y la dificultad para iniciarlos. En el contexto del diseño de software para \textit{e-Health}, este síntoma constituye una restricción crítica que obliga a reconsiderar los tiempos de respuesta de la interfaz. Teóricamente, implica que el sistema debe evitar tiempos de espera cortos y eliminar gestos complejos o rápidos, favoreciendo interacciones pausadas y áreas de contacto extendidas para mitigar la frustración del usuario.  \cite{parkinson_bradicinesia}

\textbf{Telerrehabilitación.} Es la rama de la telemedicina que aplica las Tecnologías de la Información y la Comunicación (TIC) para la provisión de servicios de rehabilitación a distancia. Este modelo teórico busca descentralizar la atención sanitaria, permitiendo la evaluación, el diagnóstico y el tratamiento de pacientes fuera del entorno hospitalario tradicional. Su objetivo es garantizar la continuidad asistencial, aumentar la adherencia al tratamiento mediante la monitorización remota y democratizar el acceso a terapias especializadas para pacientes con movilidad reducida o residentes en zonas geográficas aisladas.

\textbf{\textit{e-Health} o Salud Digital.} Se define como la intersección multidisciplinar entre la informática médica, la salud pública y las ciencias de la empresa referida a los servicios sanitarios y la información transmitida a través de Internet y las tecnologías relacionadas. Más allá de una simple implementación técnica, este concepto teórico representa un cambio de paradigma en el modelo asistencial que busca mejorar el acceso, la eficiencia y la calidad de la atención médica mediante la digitalización de procesos. Su aplicación abarca desde la gestión administrativa de historias clínicas electrónicas hasta la provisión de servicios de telemedicina, permitiendo la deslocalización de la asistencia sanitaria.  \cite{uoc_ehealth}

\textbf{\textit{m-Health} o Salud Móvil.} Constituye un subsegmento específico dentro de la \textit{e-Health} que se centra en la práctica de la medicina y la salud pública soportada por dispositivos móviles como teléfonos inteligentes, tabletas y tecnologías vestibles. La característica teórica distintiva de este modelo es la ubicuidad y la capacidad de aprovechar los sensores integrados en el hardware del dispositivo, como acelerómetros, giroscopios o cámaras de alta resolución, para la recolección de datos clínicos en tiempo real. En el contexto de la rehabilitación neurológica, la \textit{m-Health} democratiza el acceso al tratamiento al convertir un dispositivo de consumo masivo en una herramienta de monitorización médica continua y accesible desde el domicilio del paciente.  \cite{wiki_mhealth}

\section{Ingeniería del software y metodología}

\textbf{Metodología ágil de gestión.} Se define como una filosofía de gestión de proyectos que prioriza la flexibilidad, la colaboración y la entrega continua de valor frente a la planificación rígida a largo plazo. En lugar de abordar el desarrollo como un bloque monolítico, se estructura el trabajo en ciclos temporales cortos y repetitivos. Este enfoque permite la inspección constante del progreso y la adaptación rápida a los cambios en los requisitos, asegurando que el producto final se alinee con las necesidades reales del usuario.

\textbf{\textit{Framework}.} En la ingeniería de software, se define como una estructura conceptual y tecnológica de soporte que sirve de base para la organización y desarrollo de proyectos complejos. A diferencia de una simple librería de código, el \textit{framework} se caracteriza por imponer el principio de inversión de control. Esto significa que no es el código del programador el que llama a las funciones de la herramienta, sino que es el propio entorno el que orquesta el flujo de ejecución general y delega en el código personalizado del desarrollador únicamente los detalles específicos de la lógica de negocio. Su adopción permite estandarizar la arquitectura del software, reutilizar componentes probados y aumentar la productividad al abstraer tareas repetitivas de bajo nivel.  \cite{unir_framework}

\textbf{Desarrollo \textit{frontend}.} Hace referencia a la ingeniería de la capa de presentación del software. Comprende la implementación de la interfaz gráfica y la lógica de interacción que se ejecuta directamente en el dispositivo del usuario final. Su objetivo teórico es garantizar la usabilidad y la accesibilidad, transformando los datos recibidos del servidor en una experiencia visual comprensible y reactiva con la que el usuario puede interactuar.

\textbf{Desarrollo \textit{backend}.} Constituye la ingeniería de la capa de acceso a datos y lógica de negocio que opera en un servidor remoto, oculta al usuario. Se encarga de procesar las reglas del sistema, gestionar la seguridad, realizar cálculos complejos y asegurar la persistencia de la información en la base de datos. Es el responsable de centralizar la verdad del sistema y servir los datos solicitados por las aplicaciones cliente de manera segura y eficiente.

\section{Patrones de diseño y arquitectura aplicada}

\textbf{Patrón de arquitectura de software.} Este concepto hace referencia a las reglas organizativas que definen la estructura interna del código fuente. Su objetivo fundamental es la separación de responsabilidades, buscando desacoplar la lógica de negocio y la gestión de datos de la interfaz gráfica. Al dividir el sistema en capas independientes con funciones específicas, se facilita el mantenimiento del código, la escalabilidad del proyecto y la realización de pruebas unitarias aisladas.

\textbf{Patrón de arquitectura \textit{Model-View-ViewModel} (MVVM).} Es un patrón de diseño arquitectónico cuyo objetivo fundamental es desacoplar la lógica de presentación y el estado de la interfaz de usuario de la lógica de negocio y los datos. Se estructura en tres componentes diferenciados: el Modelo, que representa los datos y las reglas de dominio; la Vista, que es un componente pasivo encargado de la representación visual; y el \textit{ViewModel}, que actúa como intermediario gestionando el estado de la aplicación y procesando las interacciones. Su característica teórica distintiva es la ausencia de referencia directa desde el \textit{ViewModel} hacia la Vista, basándose en mecanismos de observación para actualizar la interfaz automáticamente ante cambios en los datos.

\textbf{Patrón de diseño \textit{Data Access Object} (DAO).} Es un patrón estructural utilizado en ingeniería de software para aislar la lógica de negocio de la capa de persistencia. Su propósito teórico es abstraer y encapsular todos los accesos a la fuente de datos. El DAO gestiona la conexión técnica con el mecanismo de almacenamiento (ya sea una base de datos relacional, un archivo local o una API externa) para realizar operaciones de lectura y escritura, exponiendo una interfaz pública agnóstica a la implementación subyacente. Esto permite que el resto de la aplicación opere independientemente de la tecnología de base de datos específica, facilitando la mantenibilidad y la migración de sistemas.

\textbf{Paradigma de programación declarativa.} Se trata de un modelo de programación enfocado en el resultado final visual en lugar de en el flujo de control paso a paso. En este paradigma, el desarrollador describe qué elementos deben aparecer en la pantalla en función del estado actual de la aplicación, delegando en el sistema la responsabilidad de cómo renderizar o actualizar dichos elementos. La interfaz gráfica se concibe como una función directa e inmutable de los datos, actualizándose automáticamente cuando estos varían.

\section{Conceptos Técnicos de Implementación}

\textbf{Ciclo de vida de la aplicación.} Comprende la secuencia de estados por los que transita un proceso de software dentro del sistema operativo, desde su creación hasta su destrucción. Este concepto teórico dicta que una aplicación no es una entidad estática, sino que reacciona a cambios del entorno, como pasar a segundo plano o ser interrumpida. Una correcta gestión de estos estados es crucial para liberar recursos del sistema, como la cámara o la memoria, cuando la aplicación deja de ser visible para el usuario.

\textbf{Concurrencia estructurada y programación asíncrona.} Este concepto computacional se refiere a la capacidad del software para ejecutar tareas múltiples de manera paralela sin bloquear el hilo principal de ejecución. Es fundamental en entornos con interfaz gráfica para gestionar operaciones de larga duración, como el acceso a redes o el procesamiento de archivos, en un segundo plano. La concurrencia estructurada organiza estas tareas jerárquicamente, garantizando que los procesos secundarios se completen o cancelen de forma ordenada para evitar fugas de memoria.

\textbf{\textit{Endpoint}.} Constituye el extremo final de un canal de comunicación dentro de una Interfaz de Programación de Aplicaciones (API). En el contexto de los servicios web, un \textit{endpoint} se materializa como una dirección URL específica unida a un método de petición concreto, como GET o POST, que permite a los sistemas clientes interactuar con un recurso determinado del servidor. Desde una perspectiva teórica, representa la interfaz de contrato pública del sistema, definiendo estrictamente la ubicación donde los recursos están disponibles y el formato en el que se debe intercambiar la información, actuando como la puerta de entrada controlada y segura hacia la lógica interna de la aplicación.

\textbf{Dibujo vectorial.} Consiste en una técnica de representación gráfica basada en fórmulas matemáticas y primitivas geométricas, como líneas, curvas y polígonos, en lugar de utilizar matrices de píxeles estáticas. La principal ventaja teórica de este modelo es la independencia de la resolución, lo que permite que los elementos visuales se redimensionen y escalen a cualquier tamaño de pantalla sin sufrir pérdida de nitidez o pixelación.

\section{Seguridad y privacidad}

\textbf{Aislamiento de aplicaciones o \textit{Sandboxing}.} Es un principio fundamental de seguridad informática en sistemas operativos modernos. Establece que cada aplicación se ejecuta en un entorno restringido y privado, con su propio espacio de almacenamiento y memoria asignados. Bajo este modelo teórico, ningún proceso puede acceder a los datos o archivos generados por otra aplicación sin una autorización explícita, garantizando la privacidad y la integridad de la información del usuario.  \cite{esed_sandboxing}

\textbf{Cifrado de información.} Es la disciplina teórica y matemática responsable de la seguridad de la información mediante la ofuscación de datos. Se define como el procedimiento algorítmico mediante el cual un mensaje legible o texto plano se transforma en un conjunto de caracteres aparentemente aleatorios o texto cifrado. Este proceso garantiza el principio de confidencialidad, asegurando que la información solo pueda ser revertida a su estado original por entidades que posean la clave criptográfica correspondiente, protegiendo así los datos frente a accesos no autorizados durante su almacenamiento o transmisión.

\textbf{\textit{Hashing} de cifrado.} Hace referencia a un tipo específico de algoritmo matemático que transforma cualquier bloque de datos arbitrario en una cadena de caracteres de longitud fija y estructura única, denominada \textit{hash} o huella digital. A diferencia del cifrado convencional, su característica teórica fundamental es la unidireccionalidad o irreversibilidad: es computacionalmente inviable recuperar los datos originales a partir del \textit{hash} resultante. Este concepto es el estándar de la industria para el almacenamiento seguro de credenciales y contraseñas, ya que permite verificar la autenticidad de un usuario sin necesidad de guardar su clave real en la base de datos. 